% GENERATED FILE. Edit canonical lessons, then run npm run generate:books.
% canonical-source-hash: fnv1a64:e8a1d03b7761eec6
% canonical-lessons: BN-C412-pachil, BN-C412-get, BN-C412-dalan, BN-C412-lemppost, BN-C412-topkano

\chapter{Walls and Gates}
\label{ch:bn-walls-and-gates}
\hlchaptermodality{412}

\begin{chapteropening}
\textbf{By the end of this chapter:} \emph{I can describe the walls, gates and buildings of a street.}

Complete the last lesson of chapter 412: I can describe the walls, gates and buildings of a street.
\end{chapteropening}

\section[\texorpdfstring{\bn{পাঁচিল} (pā̃chil)}{পাঁচিল (pā̃chil)}]{\bn{পাঁচিল} (pā̃chil) — a boundary wall}

\label{lesson:BN-C412-pachil}



\begin{quote}
Before the new one: say the Bengali for a postman, then the Bengali for to send.
\end{quote}

\subsection*{You'll want to know: \bn{পাঁচিল}}
\textbf{\bn{পাঁচিল}} — \emph{pā̃chil} — \textquotedblleft{}a boundary wall\textquotedblright{}.

\begin{quote}
\bn{বাগানের} \bn{চারপাশে} \bn{উঁচু} \bn{পাঁচিল।}
\end{quote}

\begin{quote}
\emph{bāgāner chārpāshe ũchu pā̃chil.}
\end{quote}

\begin{quote}
There is a high wall all round the garden.
\end{quote}

The outer wall round a yard or garden; \bn{দেয়াল} is any wall, inside or out.

\begin{cousinweb}
\textbf{\bn{পাঁচিল}} (\emph{pā̃chil}) is inherited: CDIAL lists Bengali \emph{pā̃cil}, \textquotedblleft{}a rampart\textquotedblright{}, under Sanskrit \emph{prācīra}, \textquotedblleft{}an enclosure, a wall\textquotedblright{} (CDIAL 8921); \bn{প্রাচীর}, the same word taken whole, is the grander city wall.
\end{cousinweb}

\begin{grammarlens}[title={What you've built}]
One more word for getting around.
\end{grammarlens}

\subsection*{Your turn}
\begin{itemize}
\raggedright
  \item \emph{Say it:} \emph{pā̃chil}
  \item \emph{Say it:} \emph{pā̃chil}, once more
  \item \emph{Say it:} \emph{bāgāner chārpāshe ũchu pā̃chil.} — There is a high wall all round the garden.
  \item \emph{Recall:} say the Bengali for a postman, then the Bengali for to send, then say \emph{pā̃chil} again
\end{itemize}

\begin{tcolorbox}[breakable,colback=teal!4,colframe=teal!35!black,title={Before you move on}]
What does \bn{পাঁচিল} mean? (\textquotedblleft{}A boundary wall\textquotedblright{}.) Say it once more.
\end{tcolorbox}
\section[\texorpdfstring{\bn{গেট} (geṭ)}{গেট (geṭ)}]{\bn{গেট} (geṭ) — a gate}

\label{lesson:BN-C412-get}



\begin{quote}
Before the new one: say the Bengali for to send, then the Bengali for a boundary wall.
\end{quote}

\subsection*{You'll want to know: \bn{গেট}}
\textbf{\bn{গেট}} — \emph{geṭ} — \textquotedblleft{}a gate\textquotedblright{}.

\begin{quote}
\bn{দারোয়ান} \bn{রাতে} \bn{গেট} \bn{বন্ধ} \bn{করেন।}
\end{quote}

\begin{quote}
\emph{dāroyān rāte geṭ bôndhô kôren.}
\end{quote}

\begin{quote}
The gatekeeper shuts the gate at night.
\end{quote}

\begin{cousinweb}
\textbf{\bn{গেট}} (\emph{geṭ}) is borrowed from English \emph{gate}, with the final t said as \bn{ট}.
\end{cousinweb}

\begin{grammarlens}[title={What you've built}]
One more word for getting around.
\end{grammarlens}

\subsection*{Your turn}
\begin{itemize}
\raggedright
  \item \emph{Say it:} \emph{geṭ}
  \item \emph{Say it:} \emph{geṭ}, once more
  \item \emph{Say it:} \emph{dāroyān rāte geṭ bôndhô kôren.} — The gatekeeper shuts the gate at night.
  \item \emph{Recall:} say the Bengali for to send, then the Bengali for a boundary wall, then say \emph{geṭ} again
\end{itemize}

\begin{tcolorbox}[breakable,colback=teal!4,colframe=teal!35!black,title={Before you move on}]
What does \bn{গেট} mean? (\textquotedblleft{}A gate\textquotedblright{}.) Say it once more.
\end{tcolorbox}
\section[\texorpdfstring{\bn{দালান} (dālān)}{দালান (dālān)}]{\bn{দালান} (dālān) — a brick building; a hall}

\label{lesson:BN-C412-dalan}



\begin{quote}
Before the new one: say the Bengali for a boundary wall, then the Bengali for a gate.
\end{quote}

\subsection*{You'll want to know: \bn{দালান}}
\textbf{\bn{দালান}} — \emph{dālān} — \textquotedblleft{}a brick building; a hall\textquotedblright{}.

\begin{quote}
\bn{রাস্তার} \bn{পাশে} \bn{একটা} \bn{পুরনো} \bn{দালান} \bn{আছে।}
\end{quote}

\begin{quote}
\emph{rāstār pāshe ekṭā purôno dālān āchhe.}
\end{quote}

\begin{quote}
There is an old brick building by the road.
\end{quote}

In a village, the one house built of brick rather than mud and thatch.

\begin{cousinweb}
\textbf{\bn{দালান}} (\emph{dālān}) is usually traced to Persian \emph{dālān}, \textquotedblleft{}a hall, a covered porch\textquotedblright{}.
\end{cousinweb}

\begin{grammarlens}[title={What you've built}]
One more word for getting around.
\end{grammarlens}

\subsection*{Your turn}
\begin{itemize}
\raggedright
  \item \emph{Say it:} \emph{dālān}
  \item \emph{Say it:} \emph{dālān}, once more
  \item \emph{Say it:} \emph{rāstār pāshe ekṭā purôno dālān āchhe.} — There is an old brick building by the road.
  \item \emph{Recall:} say the Bengali for a boundary wall, then the Bengali for a gate, then say \emph{dālān} again
\end{itemize}

\begin{tcolorbox}[breakable,colback=teal!4,colframe=teal!35!black,title={Before you move on}]
What does \bn{দালান} mean? (\textquotedblleft{}A brick building; a hall\textquotedblright{}.) Say it once more.
\end{tcolorbox}
\section[\texorpdfstring{\bn{ল্যাম্পপোস্ট} (læmpposṭ)}{ল্যাম্পপোস্ট (læmpposṭ)}]{\bn{ল্যাম্পপোস্ট} (læmpposṭ) — a lamp post}

\label{lesson:BN-C412-lemppost}



\begin{quote}
Before the new one: say the Bengali for a gate, then the Bengali for a brick building.
\end{quote}

\subsection*{You'll want to know: \bn{ল্যাম্পপোস্ট}}
\textbf{\bn{ল্যাম্পপোস্ট}} — \emph{læmpposṭ} — \textquotedblleft{}a lamp post\textquotedblright{}.

\begin{quote}
\bn{গলির} \bn{মোড়ে} \bn{একটা} \bn{ল্যাম্পপোস্ট} \bn{আছে।}
\end{quote}

\begin{quote}
\emph{gôlir moṛe ekṭā læmpposṭ āchhe.}
\end{quote}

\begin{quote}
There is a lamp post at the corner of the lane.
\end{quote}

\begin{cousinweb}
\textbf{\bn{ল্যাম্পপোস্ট}} (\emph{læmpposṭ}) is borrowed from English \emph{lamp post}, written as one word.
\end{cousinweb}

\begin{grammarlens}[title={What you've built}]
One more word for getting around.
\end{grammarlens}

\subsection*{Your turn}
\begin{itemize}
\raggedright
  \item \emph{Say it:} \emph{læmpposṭ}
  \item \emph{Say it:} \emph{læmpposṭ}, once more
  \item \emph{Say it:} \emph{gôlir moṛe ekṭā læmpposṭ āchhe.} — There is a lamp post at the corner of the lane.
  \item \emph{Recall:} say the Bengali for a gate, then the Bengali for a brick building, then say \emph{læmpposṭ} again
\end{itemize}

\begin{tcolorbox}[breakable,colback=teal!4,colframe=teal!35!black,title={Before you move on}]
What does \bn{ল্যাম্পপোস্ট} mean? (\textquotedblleft{}A lamp post\textquotedblright{}.) Say it once more.
\end{tcolorbox}
\section[\texorpdfstring{\bn{টপকানো} (ṭôpkāno)}{টপকানো (ṭôpkāno)}]{\bn{টপকানো} (ṭôpkāno) — to climb over; to jump over}

\label{lesson:BN-C412-topkano}



\begin{quote}
Before the new one: say the Bengali for a brick building, then the Bengali for a lamp post.
\end{quote}

\subsection*{You'll want to know: \bn{টপকানো}}
\textbf{\bn{টপকানো}} — \emph{ṭôpkāno} — \textquotedblleft{}to climb over; to jump over\textquotedblright{}.

\begin{quote}
\bn{চোর} \bn{পাঁচিল} \bn{টপকে} \bn{পালাল।}
\end{quote}

\begin{quote}
\emph{chor pā̃chil ṭôpke pālālô.}
\end{quote}

\begin{quote}
The thief climbed over the wall and ran away.
\end{quote}

Past (he/she …): \bn{তিনি} \bn{পাঁচিল} \bn{টপকালেন}, \emph{tini pā̃chil ṭôpkālen}.

\begin{cousinweb}
\textbf{\bn{টপকানো}} (\emph{ṭôpkāno}) is inherited: CDIAL lists Bengali \emph{ṭapkāna}, \textquotedblleft{}to leap over\textquotedblright{}, under the reconstructed \emph{tarpati}, \textquotedblleft{}jumps\textquotedblright{} (CDIAL 5727).
\end{cousinweb}

\begin{grammarlens}[title={What you've built}]
One more word for getting around.
\end{grammarlens}

\subsection*{Your turn}
\begin{itemize}
\raggedright
  \item \emph{Say it:} \emph{ṭôpkāno}
  \item \emph{Say it:} \emph{ṭôpkāno}, once more
  \item \emph{Say it:} \emph{chor pā̃chil ṭôpke pālālô.} — The thief climbed over the wall and ran away.
  \item \emph{Recall:} say the Bengali for a brick building, then the Bengali for a lamp post, then say \emph{ṭôpkāno} again
\end{itemize}

\begin{tcolorbox}[breakable,colback=teal!4,colframe=teal!35!black,title={Before you move on}]
What does \bn{টপকানো} mean? (\textquotedblleft{}To climb over; to jump over\textquotedblright{}.) Say it once more.
\end{tcolorbox}
